% ---------------------------------------------------------------------------
% Sprint 1 — Documento de Visão e Definição da Solução
% Sentinela Nuclear Challenge (AMAZUL) — Grupo NextGen — 2CCPS — FIAP
% Formatação segundo ABNT NBR 14724 (estrutura) e NBR 6023 (referências)
% ---------------------------------------------------------------------------
\documentclass[12pt,a4paper,oneside]{article}

\usepackage[T1]{fontenc}
\usepackage[utf8]{inputenc}
\usepackage{times}
\usepackage[top=3cm,bottom=2cm,left=3cm,right=2cm]{geometry}
\usepackage{setspace}
\usepackage{indentfirst}
\usepackage{titlesec}
\usepackage{tocloft}
\usepackage{graphicx}
\usepackage{float}
\usepackage{xcolor}
\usepackage{array}
\usepackage{tabularx}
\usepackage{booktabs}
\usepackage{longtable}
\usepackage{enumitem}
\usepackage{caption}
\usepackage{fancyhdr}
\usepackage{amssymb}
\usepackage{hyperref}
\hypersetup{hidelinks}

\renewcommand{\contentsname}{SUMÁRIO}
\renewcommand{\figurename}{Figura}
\renewcommand{\tablename}{Tabela}
\hyphenpenalty=4000
\tolerance=2500
\emergencystretch=3em

\titleformat{\section}{\normalfont\bfseries\large}{\thesection}{1em}{}
\titleformat{\subsection}{\normalfont\bfseries\normalsize}{\thesubsection}{1em}{}
\titleformat{\subsubsection}{\normalfont\normalsize\itshape}{\thesubsubsection}{1em}{}
\titlespacing*{\section}{0pt}{1.5\baselineskip}{1\baselineskip}
\titlespacing*{\subsection}{0pt}{1.2\baselineskip}{0.8\baselineskip}
\titlespacing*{\subsubsection}{0pt}{1\baselineskip}{0.6\baselineskip}

\renewcommand{\cfttoctitlefont}{\hfill\bfseries\large}
\renewcommand{\cftaftertoctitle}{\hfill}
\setlength{\cftbeforesecskip}{4pt}
\renewcommand{\cftsecfont}{\normalfont}
\renewcommand{\cftsecpagefont}{\normalfont}
\renewcommand{\cftsecleader}{\cftdotfill{\cftdotsep}}

\captionsetup{labelsep=endash,font=small,justification=centering,
              singlelinecheck=false,skip=4pt}
\newcommand{\fonte}[1]{\par\vspace{5pt}{\footnotesize\singlespacing Fonte: #1\par}}

\pagestyle{fancy}
\fancyhf{}
\fancyhead[R]{\small\thepage}
\renewcommand{\headrulewidth}{0pt}
\setlength{\headheight}{14pt}

\newcolumntype{L}[1]{>{\raggedright\arraybackslash}p{#1}}
\newcolumntype{C}[1]{>{\centering\arraybackslash}p{#1}}
\renewcommand{\tabularxcolumn}[1]{>{\raggedright\arraybackslash}p{#1}}
\newcommand{\cab}[1]{\textbf{\footnotesize #1}}

\begin{document}
\onehalfspacing

% =========================================================== CAPA
\thispagestyle{empty}
\begin{center}
{\large\bfseries FACULDADE DE INFORMÁTICA E ADMINISTRAÇÃO PAULISTA — FIAP}\\[2pt]
{\normalsize CIÊNCIA DA COMPUTAÇÃO — TURMA 2CCPS}\\[2pt]
{\normalsize SOFTWARE ENGINEERING — GRUPO NEXTGEN}

\vspace{1.4cm}
{\normalsize
GUSTAVO FRANZOTI GONÇALVES — RM566983\\
LINCOLN SIMÃO PEREIRA — RM567284\\
MAYKON SANTANA FONSECA — RM567041\\
NICOLAS SAKAUE NISHIMURA — RM567752\par}

\vspace{2.6cm}
{\LARGE\bfseries ARGOS\par}
\vspace{0.4cm}
{\large Análise de Risco e Guarda Operacional por Sensoriamento\par}
\vspace{0.8cm}
{\normalsize\bfseries DOCUMENTO DE VISÃO E DEFINIÇÃO DA SOLUÇÃO\par}
\vspace{0.3cm}
{\normalsize Sprint 1 — Sentinela Nuclear Challenge\par}
\vspace{0.2cm}
{\normalsize AMAZUL — Amazônia Azul Tecnologias de Defesa S.A.\par}

\vfill
{\normalsize SÃO PAULO\\2026\par}
\end{center}
\newpage

% =========================================================== FOLHA DE ROSTO
\thispagestyle{empty}
\begin{center}
{\normalsize GRUPO NEXTGEN — 2CCPS\par}
\vspace{2.4cm}
{\large\bfseries ARGOS: DOCUMENTO DE VISÃO E DEFINIÇÃO DA SOLUÇÃO\par}
\vspace{0.4cm}
{\normalsize Sprint 1 do Sentinela Nuclear Challenge\par}
\end{center}

\vspace{1.4cm}
\begin{flushright}
\begin{minipage}{0.54\textwidth}
\singlespacing\small
Documento de Visão elaborado para a Sprint 1 da disciplina de Software
Engineering, do curso de Ciência da Computação da FIAP, em resposta ao desafio
Sentinela Nuclear proposto pela AMAZUL — Amazônia Azul Tecnologias de Defesa
S.A.

\vspace{6pt}
Prof.\ Leandro Ramos da Silva
\end{minipage}
\end{flushright}

\vspace{1.4cm}
\begin{center}
\begin{minipage}{0.88\textwidth}
\singlespacing\small
\textbf{Integrantes do grupo}\\[6pt]
\begin{tabularx}{\textwidth}{L{6.6cm} L{3.0cm} X}
\toprule
\cab{Nome} & \cab{RM} & \cab{Atribuição nesta Sprint} \\
\midrule
Gustavo Franzoti Gonçalves & RM566983 & Análise do contexto e mapeamento de
stakeholders \\
Lincoln Simão Pereira & RM567284 & Definição da solução, arquitetura
conceitual e repositório \\
Maykon Santana Fonseca & RM567041 & Escopo, requisitos iniciais e matriz de
rastreabilidade \\
Nicolas Sakaue Nishimura & RM567752 & Protótipo conceitual e fluxo de
utilização \\
\bottomrule
\end{tabularx}

\vspace{12pt}
\textbf{Repositório do projeto:}
\texttt{https://github.com/lincoln743/sprint01\_SE}\\[4pt]
\textbf{Versão deste documento:} 1.0 — 30 de outubro de 2026
\end{minipage}
\end{center}

\vfill
\begin{center}
{\normalsize SÃO PAULO\\2026\par}
\end{center}
\newpage

% =========================================================== SUMÁRIO
\thispagestyle{empty}
\tableofcontents
\newpage
\setcounter{page}{4}

% =========================================================== 1
\section{INTRODUÇÃO}

\subsection{Apresentação}

A AMAZUL — Amazônia Azul Tecnologias de Defesa S.A. é uma empresa pública
vinculada ao Ministério da Defesa por meio do Comando da Marinha, criada pela
Lei n.º 12.706, de 8 de agosto de 2012 (BRASIL, 2012), para desenvolver,
transferir e manter tecnologias de interesse estratégico nacional —
especialmente as ligadas ao Programa Nuclear da Marinha e ao Programa de
Desenvolvimento de Submarinos. Instalações dessa natureza concentram material,
conhecimento e infraestrutura cujo comprometimento teria consequências que
excedem em muito o prejuízo patrimonial.

O Sentinela Nuclear Challenge propõe a seguinte pergunta: como utilizar drones,
sensores, visão computacional e Inteligência Artificial para monitorar,
inspecionar e identificar situações anômalas em ambientes críticos relacionados
à segurança física de instalações nucleares?

Este documento é a resposta do Grupo NextGen à Sprint 1 do desafio. Ele não
trata de implementação: trata de entender o problema, identificar quem é
afetado por ele, propor uma solução coerente e delimitar com honestidade o que
essa solução fará e o que não fará.

\subsection{Objetivo do documento}

O objetivo é apresentar a visão inicial do \textbf{ARGOS} — Análise de Risco e
Guarda Operacional por Sensoriamento —, registrando o problema a ser resolvido,
os usuários envolvidos, as necessidades identificadas, a proposta de solução da
equipe, os benefícios esperados e o escopo inicial do sistema, de modo que o
documento sirva de referência estável para as etapas seguintes de análise,
modelagem e desenvolvimento.

A escolha do nome tem uma razão: na mitologia grega, Argos Panoptes era o
guardião de cem olhos, dos quais apenas alguns dormiam por vez. É uma boa
descrição do que se pretende — vigilância contínua, distribuída e sem ponto
cego. O acrônimo em português descreve as três funções do sistema: análise de
risco, guarda operacional e sensoriamento.

\subsection{Organização do documento}

A Seção~\ref{sec:contexto} analisa o contexto e o problema. A
Seção~\ref{sec:objetivos} declara os objetivos e os indicadores pelos quais a
solução deverá ser julgada. A Seção~\ref{sec:stakeholders} mapeia os
envolvidos. A Seção~\ref{sec:solucao} descreve a solução proposta, sua
arquitetura conceitual, suas funcionalidades, as tecnologias previstas e os
diferenciais. A Seção~\ref{sec:beneficios} trata dos benefícios esperados e a
Seção~\ref{sec:escopo}, do escopo — inclusive do que está deliberadamente fora
dele. A Seção~\ref{sec:prototipo} apresenta o protótipo conceitual de baixa
fidelidade e a Seção~\ref{sec:conclusao} encerra com os próximos passos. O
Apêndice~A traz a matriz de rastreabilidade entre necessidades, requisitos e
benefícios.

\subsection{Método de trabalho}

A equipe seguiu, nesta Sprint, o roteiro proposto na disciplina: compreensão do
problema a partir do enunciado e da pesquisa documental sobre o domínio;
identificação dos stakeholders e de suas necessidades; definição da proposta de
valor; delimitação do escopo; e prototipação conceitual. As decisões de maior
consequência foram registradas no formato de ADR (\emph{Architecture Decision
Record}) proposto por Nygard (2011), prática já adotada pelo grupo nos
artefatos anteriores da disciplina — a Seção~\ref{sec:adr} traz o primeiro
registro.

Vale um esclarecimento metodológico. O grupo não teve acesso a dados
operacionais reais de nenhuma instalação da AMAZUL, por razões evidentes de
segurança. Os números apresentados como linha de base são \emph{estimativas da
equipe}, construídas a partir de literatura pública e do enunciado do desafio,
e estão sempre identificados como tal. Nenhum deles deve ser lido como medição.

% =========================================================== 2
\section{CONTEXTO DO PROBLEMA}
\label{sec:contexto}

\subsection{O que é segurança física em uma instalação nuclear}

Convém separar dois conceitos que o vocabulário cotidiano confunde. A
\emph{segurança nuclear} (\emph{safety}) trata de evitar que o processo falhe e
cause dano — é o domínio dos sistemas de proteção do reator. A \emph{proteção
física} (\emph{security}) trata de evitar que alguém, deliberadamente, cause
dano — é o domínio de cercas, controle de acesso, vigilância e resposta. O
desafio proposto pela AMAZUL situa-se inteiramente no segundo campo.

A norma CNEN NN 2.01 — Proteção Física de Materiais e Instalações Nucleares,
aprovada pela Resolução CNEN n.º 253/2019 (COMISSÃO NACIONAL DE ENERGIA
NUCLEAR, 2019) — estabelece os princípios e requisitos para o projeto de um
Sistema de Proteção Física no Brasil. No plano internacional, a recomendação de
referência é a INFCIRC/225/Revisão 5, publicada pela Agência Internacional de
Energia Atômica como Nuclear Security Series n.º 13 (INTERNATIONAL ATOMIC
ENERGY AGENCY, 2011). Ambas organizam a proteção física em torno de três
funções encadeadas, que estruturam todo o raciocínio deste documento:

\begin{enumerate}[leftmargin=*,itemsep=2pt,topsep=4pt]
  \item \textbf{detecção} — perceber que algo anômalo está acontecendo e
        avaliar o que é;
  \item \textbf{atraso} — barreiras físicas que aumentam o tempo de que o
        adversário precisa;
  \item \textbf{resposta} — interromper a ação dentro do tempo que o atraso
        concedeu.
\end{enumerate}

A relação entre as três é aritmética e implacável: a proteção só é efetiva se o
tempo de detecção somado ao tempo de resposta for menor que o tempo de atraso.
Barreiras físicas são caras e lentas de modificar; o tempo de resposta depende
de efetivo e de treinamento. Sobra a detecção — e é exatamente nela que a
tecnologia proposta pelo desafio pode produzir ganho mensurável e rápido.
\textbf{Este é o problema que o ARGOS se propõe a atacar: reduzir o tempo entre
o evento e a sua compreensão pelo operador.}

\subsection{Riscos e situações anômalas consideradas}

A Tabela~\ref{tab:anomalias} relaciona as situações anômalas que o grupo
considerou plausíveis no recorte de segurança física. A coluna de detecção
atual descreve o cenário típico de uma instalação equipada com circuito fechado
de televisão convencional e rondas periódicas.

\begin{table}[H]
\centering
\caption{Situações anômalas consideradas no recorte do desafio}
\label{tab:anomalias}
{\small\singlespacing
\begin{tabularx}{\textwidth}{L{3.3cm} X L{3.5cm}}
\toprule
\cab{Anomalia} & \cab{Descrição e consequência potencial} & \cab{Como é
detectada hoje} \\
\midrule
Intrusão no perímetro & Pessoa cruza ou tenta cruzar a cerca em ponto não
vigiado; é o evento inicial de praticamente todo cenário de ameaça &
Sensor de cerca com alta taxa de alarme falso; confirmação visual depende de o
operador estar olhando a câmera certa \\
\addlinespace
Permanência em área restrita & Pessoa credenciada para uma área é detectada em
outra, ou permanece além do tempo previsto & Controle de acesso registra a
entrada, mas não acompanha a permanência \\
\addlinespace
Sobrevoo não autorizado & Aeronave não tripulada sobre a instalação, para
reconhecimento ou transporte de carga & Praticamente não detectada; depende de
observação casual \\
\addlinespace
Veículo em comportamento atípico & Veículo parado junto ao perímetro, em
circulação fora de rota ou fora de horário & Depende da ronda ou de o operador
notar a imagem \\
\addlinespace
Objeto abandonado & Volume deixado junto a acesso, galpão ou cerca & Depende de
inspeção visual presencial \\
\addlinespace
Violação de barreira ou lacre & Porta, portão, grade ou lacre violado,
inclusive por ação interna & Verificada em ronda, com atraso de horas \\
\addlinespace
Degradação do próprio sistema & Câmera cega, sensor mudo, enlace caído,
iluminação queimada — a falha silenciosa que anula a detecção & Percebida por
acaso, às vezes apenas quando se busca a gravação \\
\bottomrule
\end{tabularx}}
\fonte{elaborado pelos autores (2026), com base nas funções de proteção física
descritas na CNEN NN 2.01 (2019) e na INFCIRC/225/Rev.5 (2011).}
\end{table}

\subsection{Por que o monitoramento convencional não basta}

Três limitações do arranjo convencional explicam por que acrescentar câmeras,
isoladamente, não resolve o problema.

A primeira é a atenção humana. A literatura de fatores humanos em vigilância é
consistente ao apontar queda significativa de desempenho na detecção de eventos
raros após os primeiros vinte a trinta minutos de observação contínua — e a
operação de monitoramento é, por definição, a espera de um evento raro. Um
operador responsável por dezenas de câmeras simultâneas não falha por
descuido; falha porque a tarefa, como desenhada, excede a capacidade atencional
humana.

A segunda é o volume. Cada câmera acrescentada melhora a cobertura e piora a
atenção disponível por câmera. Sem triagem automática, ampliar o parque de
sensores aumenta a quantidade de informação não observada.

A terceira é a natureza do registro. O CFTV convencional é essencialmente
\emph{forense}: serve para reconstituir o que aconteceu depois que aconteceu.
A proteção física precisa do oposto — precisa de detecção em tempo útil para
que a resposta ocorra dentro do tempo de atraso que as barreiras oferecem.

\subsection{Como as tecnologias emergentes mudam esse quadro}

As tecnologias apontadas pelo desafio atacam diretamente as três limitações. A
\textbf{visão computacional} executa a triagem que o operador não consegue
sustentar, classificando continuamente cada quadro e chamando a atenção humana
apenas quando há motivo. O \textbf{processamento na borda} mantém essa triagem
em milissegundos e junto à fonte, sem depender de enlace com um centro de dados
externo — o que, além de mais rápido, é mais adequado a um ambiente em que o
tráfego de imagens para fora da instalação é indesejável. Os \textbf{sensores
IoT} cobrem o que a câmera não vê: vibração na cerca, abertura de porta,
presença em duto, condição ambiental. Os \textbf{drones autônomos} transformam
a inspeção, hoje presencial e periódica, em verificação sob demanda, levando um
sensor móvel ao ponto exato do alerta em segundos. A \textbf{Inteligência
Artificial} correlaciona essas fontes: um sensor de cerca disparando sozinho é
ruído, mas o mesmo sensor acompanhado de detecção humana no vídeo e de nenhuma
credencial registrada naquele setor é um evento de risco alto. E a
\textbf{geolocalização com painel operacional} entrega essa conclusão em forma
de ação, no mapa, com o contexto necessário para decidir.

\subsection{Impacto esperado na operação}

O impacto pretendido não é substituir o operador: é mudar a natureza do seu
trabalho. Em vez de vigiar continuamente dezenas de imagens à espera de algo,
ele passa a julgar um pequeno número de eventos já triados, com contexto
reunido e recomendação de ação. A decisão permanece humana — e, em um domínio
regulado como este, deve permanecer —, mas deixa de depender de atenção
sustentada por turnos inteiros.

% =========================================================== 3
\section{OBJETIVOS DA SOLUÇÃO}
\label{sec:objetivos}

\subsection{Objetivo geral}

Reduzir o tempo entre a ocorrência de uma situação anômala na área física de
uma instalação crítica e a sua compreensão pelo operador, por meio de um
sistema que integre sensores, câmeras com visão computacional, drones
autônomos e correlação por Inteligência Artificial em um único painel
operacional, mantendo a decisão final sob responsabilidade humana e a trilha
de cada evento auditável.

\subsection{Objetivos específicos}

\begin{enumerate}[label=\textbf{OE-\arabic*},leftmargin=*,itemsep=3pt,topsep=4pt]
  \item detectar automaticamente intrusão, permanência indevida, objeto
        abandonado, sobrevoo não autorizado e violação de barreira, a partir da
        fusão de múltiplas fontes;
  \item classificar cada detecção em um nível de severidade, de modo que a fila
        do operador seja ordenada por risco e não por ordem de chegada;
  \item reduzir o volume de alarmes falsos que chegam ao operador, aproveitando
        a correlação entre fontes independentes;
  \item permitir a verificação remota de um alerta em segundos, despachando um
        drone à coordenada em vez de deslocar uma equipe;
  \item apresentar, em um único painel, a posição geográfica, a imagem, a
        trilha do evento e as ações disponíveis;
  \item registrar toda ocorrência e toda ação do operador em trilha imutável,
        apta a compor evidência para auditoria e para o órgão regulador;
  \item monitorar a saúde do próprio sistema, tratando sensor mudo e câmera
        cega como eventos de segurança, não como mera falha de manutenção.
\end{enumerate}

\subsection{Indicadores de sucesso}

A Tabela~\ref{tab:indicadores} declara como a equipe pretende julgar a solução.
Reitere-se que a coluna de linha de base é estimativa, não medição: ela existe
para tornar a meta falsificável, e deverá ser substituída por valores reais
assim que houver acesso a dados de operação.

\begin{table}[H]
\centering
\caption{Indicadores de sucesso propostos para o ARGOS}
\label{tab:indicadores}
{\small\singlespacing
\begin{tabularx}{\textwidth}{L{4.4cm} X L{2.6cm} L{2.0cm}}
\toprule
\cab{Indicador} & \cab{O que mede} & \cab{Linha de base estimada} & \cab{Meta} \\
\midrule
Tempo até a detecção & Do instante do evento até o alerta no painel & 3 a 10
min (ronda ou acaso) & abaixo de 60 s \\
Tempo até a verificação & Do alerta até a confirmação visual do que ocorre &
10 a 20 min (deslocamento) & abaixo de 3 min \\
Taxa de alarme falso & Proporção de alertas descartados pelo operador & acima
de 50 \% (sensor isolado) & abaixo de 15 \% \\
Cobertura do perímetro & Extensão com detecção automática ativa & parcial, só
onde há câmera observada & 100 \% do perímetro \\
Disponibilidade do sistema & Tempo em que a detecção opera plenamente & não
medida & 99,5 \% mensal \\
Rastreabilidade & Ocorrências com trilha completa e exportável & registro
manual em livro & 100 \% \\
\bottomrule
\end{tabularx}}
\fonte{elaborado pelos autores (2026). Linha de base estimada pela equipe, a
ser substituída por medição.}
\end{table}

% =========================================================== 4
\section{STAKEHOLDERS}
\label{sec:stakeholders}

\subsection{Mapa dos envolvidos}

A Tabela~\ref{tab:stakeholders} relaciona os envolvidos, o papel de cada um e,
principalmente, o que cada um precisa do sistema — porque é dessa necessidade,
e não da tecnologia disponível, que os requisitos devem derivar.

\begin{table}[H]
\centering
\caption{Stakeholders do ARGOS, papéis e necessidades}
\label{tab:stakeholders}
{\small\singlespacing
\begin{tabularx}{\textwidth}{L{3.4cm} L{4.3cm} X}
\toprule
\cab{Stakeholder} & \cab{Papel} & \cab{Principal necessidade em relação ao
sistema} \\
\midrule
Operador de monitoramento & Acompanha o painel em turno contínuo e é o primeiro
a julgar cada evento & Receber poucos alertas, confiáveis e com contexto
suficiente para decidir em segundos — e não ser inundado por alarme falso \\
\addlinespace
Equipe de segurança física & Atende a ocorrência em campo & Saber exatamente
para onde ir, o que esperar encontrar e com que grau de confiança, antes de
chegar \\
\addlinespace
Equipe de inspeção e manutenção & Executa inspeções programadas e verifica
ativos & Substituir parte da inspeção presencial por inspeção aérea registrada,
com histórico comparável entre períodos \\
\addlinespace
Gestor da instalação & Responde pela segurança física perante a direção e os
órgãos externos & Indicadores confiáveis, evidência de conformidade e visão
consolidada do risco \\
\addlinespace
Coordenador do centro de operações & Dimensiona turnos e define procedimentos &
Entender a carga real de eventos para distribuir efetivo e padronizar a
resposta \\
\addlinespace
Administrador do sistema & Cadastra sensores, rotas, perfis e políticas de
severidade & Configurar sem depender do fornecedor e controlar quem vê e quem
age \\
\addlinespace
Órgãos reguladores (CNEN e, no plano internacional, AIEA) & Fiscalizam o
atendimento aos requisitos de proteção física & Registro íntegro, rastreável e
exportável das ocorrências e das respostas \\
\addlinespace
Área de TI e segurança da informação & Responde pela infraestrutura e pela
proteção dos dados & Arquitetura que não exponha imagem nem dado sensível e que
se integre ao ambiente existente \\
\addlinespace
Alta direção da AMAZUL e Comando da Marinha & Patrocina e prioriza o
investimento & Demonstração de ganho operacional e de aderência à missão
institucional \\
\addlinespace
Fornecedores de sensores e drones & Entregam e mantêm o hardware & Interfaces
abertas e documentadas, para não haver dependência de um único fabricante \\
\addlinespace
Comunidade do entorno e imprensa & Afetadas pela percepção de segurança da
instalação & Transparência sobre o que é monitorado e sobre o tratamento de
imagens \\
\bottomrule
\end{tabularx}}
\fonte{elaborado pelos autores (2026).}
\end{table}

\subsection{Priorização por influência e interesse}

Nem todos os envolvidos pesam igualmente nas decisões de projeto. A
Figura~\ref{fig:stakeholders} posiciona cada um segundo a influência que exerce
sobre o projeto e o interesse que tem no seu resultado, o que orienta a
estratégia de envolvimento de cada grupo nas próximas Sprints.

\begin{figure}[H]
\centering
\caption{Matriz de influência e interesse dos stakeholders}
\label{fig:stakeholders}
\includegraphics[width=0.96\textwidth]{fig_stakeholders.png}
\fonte{elaborado pelos autores (2026).}
\end{figure}

Duas leituras dessa matriz orientaram o restante do documento. A primeira: o
operador de monitoramento tem interesse altíssimo e influência formal baixa —
é quem mais sofrerá o efeito do sistema e quem menos participa das decisões de
compra. A equipe tratou esse desequilíbrio como risco de projeto e colocou o
operador no centro do protótipo da Seção~\ref{sec:prototipo}. A segunda: os
órgãos reguladores têm influência alta e interesse concentrado em um aspecto
específico — a evidência. Por isso a trilha de auditoria aparece como requisito
desde a Sprint 1, e não como funcionalidade a ser acrescentada no fim.

% =========================================================== 5
\section{DESCRIÇÃO DA SOLUÇÃO PROPOSTA}
\label{sec:solucao}

\subsection{Declaração de visão do produto}

\begin{quote}
\singlespacing\small
\textbf{Para} operadores de monitoramento e equipes de segurança física de
instalações críticas, \textbf{que} precisam perceber e compreender situações
anômalas em tempo útil de resposta, \textbf{o ARGOS é} uma plataforma de
vigilância integrada \textbf{que} reúne sensores, câmeras com visão
computacional e drones autônomos, correlaciona essas fontes por Inteligência
Artificial e entrega ao operador eventos já triados, classificados por risco e
georreferenciados. \textbf{Diferentemente} do circuito fechado de televisão
convencional, que registra para consulta posterior e depende da atenção
contínua de quem observa, \textbf{o ARGOS} dirige a atenção humana ao que
importa, permite verificar remotamente antes de deslocar pessoas e mantém de
cada decisão uma trilha auditável.
\end{quote}

\subsection{Arquitetura conceitual}

A Figura~\ref{fig:camadas} organiza a solução em três camadas e explicita o
percurso que o desafio enuncia: transformar dados em informação, informação em
inteligência e inteligência em ação.

\begin{figure}[H]
\centering
\caption{Arquitetura conceitual do ARGOS em três camadas}
\label{fig:camadas}
\includegraphics[width=\textwidth]{fig_camadas.png}
\fonte{elaborado pelos autores (2026).}
\end{figure}

A Figura~\ref{fig:contexto} apresenta a mesma solução pela ótica do contexto —
quem a usa e com que sistemas ela conversa. A notação segue o nível 1 do modelo
C4, adotado pelo grupo nos artefatos anteriores da disciplina; o detalhamento
em contêineres e componentes será produzido na Sprint 2, quando a arquitetura
tiver sido decidida.

\begin{figure}[H]
\centering
\caption{Diagrama de contexto do ARGOS (C4 — nível 1)}
\label{fig:contexto}
\includegraphics[width=\textwidth]{fig_contexto.png}
\fonte{elaborado pelos autores (2026).}
\end{figure}

\subsection{Principais funcionalidades}

A Tabela~\ref{tab:rf} lista as funcionalidades previstas, priorizadas pelo
método MoSCoW: \emph{essencial} para o que o sistema não pode deixar de fazer,
\emph{importante} para o que agrega valor relevante e \emph{desejável} para o
que pode ser postergado sem comprometer a proposta.

\begin{table}[H]
\centering
\caption{Funcionalidades previstas para o ARGOS}
\label{tab:rf}
{\small\singlespacing
\begin{tabularx}{\textwidth}{L{1.5cm} L{4.0cm} X L{2.2cm}}
\toprule
\cab{Código} & \cab{Funcionalidade} & \cab{Descrição} & \cab{Prioridade} \\
\midrule
RF-01 & Ingestão de sensores & Receber e normalizar eventos de sensores de
cerca, portas, movimento e condições ambientais & Essencial \\
RF-02 & Detecção por visão computacional & Classificar quadros das câmeras e
identificar pessoa, veículo e objeto abandonado & Essencial \\
RF-03 & Correlação de eventos & Cruzar fontes independentes e atribuir
pontuação de risco ao evento resultante & Essencial \\
RF-04 & Classificação de severidade & Aplicar política configurável que traduz
a pontuação em severidade alta, média ou baixa & Essencial \\
RF-05 & Painel operacional & Exibir mapa da instalação, fila de alertas e
indicadores em tempo real & Essencial \\
RF-06 & Verificação assistida & Reunir vídeo, trilha do evento, dados do
contexto e ações em uma única tela & Essencial \\
RF-07 & Registro de ocorrência & Registrar ocorrência com classificação,
evidências e ações, em trilha imutável & Essencial \\
RF-08 & Despacho de drone & Enviar drone à coordenada do alerta e transmitir
imagem ao vivo ao operador & Importante \\
RF-09 & Inspeção programada & Executar rotas periódicas de inspeção aérea e
gerar relatório comparável entre execuções & Importante \\
RF-10 & Monitoração da saúde do sistema & Detectar sensor mudo, câmera cega e
enlace caído, tratando-os como evento de segurança & Importante \\
RF-11 & Indicadores e relatórios & Consolidar indicadores operacionais e
exportar evidência para auditoria & Importante \\
RF-12 & Administração & Cadastrar sensores, rotas, perfis de acesso e políticas
de severidade & Importante \\
RF-13 & Realimentação do modelo & Registrar a classificação do operador e usá-la
para reduzir o falso positivo ao longo do tempo & Desejável \\
RF-14 & Detecção de sobrevoo não autorizado & Identificar aeronave não
tripulada no espaço aéreo da instalação & Desejável \\
\bottomrule
\end{tabularx}}
\fonte{elaborado pelos autores (2026).}
\end{table}

\subsection{Casos de uso principais}

A Figura~\ref{fig:casos} apresenta a visão preliminar dos casos de uso,
suficiente para delimitar o sistema nesta Sprint. Note-se que \emph{verificar
ocorrência com drone} e \emph{registrar e classificar a ocorrência} são
incluídos pelo caso \emph{receber e triar alerta}: a verificação e o registro
não são passos opcionais, e sim partes constitutivas do tratamento de um
alerta — decisão de projeto análoga à adotada pelo grupo no artefato de padrões
de projeto da disciplina.

\begin{figure}[H]
\centering
\caption{Casos de uso preliminares do ARGOS}
\label{fig:casos}
\includegraphics[width=\textwidth]{fig_casos_uso.png}
\fonte{elaborado pelos autores (2026).}
\end{figure}

\subsection{Tecnologias previstas}

A Tabela~\ref{tab:tecnologias} apresenta as tecnologias candidatas. Trata-se de
previsão, não de decisão: a escolha definitiva depende de avaliação técnica nas
próximas Sprints, e o documento registra o papel de cada uma para que a
discussão futura parta de um critério, e não de preferência.

\begin{table}[H]
\centering
\caption{Tecnologias candidatas por camada}
\label{tab:tecnologias}
{\small\singlespacing
\begin{tabularx}{\textwidth}{L{2.4cm} L{4.2cm} X}
\toprule
\cab{Camada} & \cab{Tecnologia candidata} & \cab{Papel e justificativa} \\
\midrule
Borda & Computador de borda com acelerador de inferência & Executar a detecção
junto à câmera, em milissegundos, sem depender de enlace externo \\
Borda & Modelos de detecção de objetos em tempo real & Identificar pessoa,
veículo e objeto; são modelos maduros, com custo computacional compatível com
a borda \\
Borda & Drones com voo autônomo por rota e câmera estabilizada & Levar um
sensor móvel ao ponto do alerta e executar inspeções programadas \\
Borda & Sensores IoT com protocolo de baixo consumo & Cobrir cerca, portas,
dutos e ambiente onde não há câmera \\
Plataforma & Fila de mensagens e processamento de eventos & Desacoplar a
captação do processamento e absorver picos sem perder evento \\
Plataforma & Mecanismo de regras e correlação com modelos de IA & Transformar
detecções isoladas em evento com pontuação de risco \\
Plataforma & Banco relacional com extensão geoespacial & Guardar ocorrências,
rotas e posições, com consulta por área \\
Plataforma & Armazenamento de objetos com retenção imutável & Preservar vídeo
e evidência sem possibilidade de alteração \\
Apresentação & Aplicação web com mapa interativo & Entregar o painel
operacional em qualquer estação do centro de operações \\
Apresentação & Serviço de notificação & Alcançar equipes em campo por
dispositivo móvel \\
Transversal & Autenticação forte com perfis e trilha de auditoria & Garantir
que cada ação tenha autor identificado e registro \\
\bottomrule
\end{tabularx}}
\fonte{elaborado pelos autores (2026).}
\end{table}

\subsection{Diferenciais da proposta}

Quatro diferenciais distinguem o ARGOS de um conjunto de câmeras com análise de
vídeo.

O primeiro é a \textbf{correlação entre fontes independentes}. Um sensor de
cerca que dispara sozinho em dia de vento é ruído; o mesmo sensor acompanhado
de detecção humana no vídeo, em setor sem credencial ativa, é evento de risco
alto. É a combinação que carrega a informação, e é por isso que o sistema é
integrado em vez de ser um conjunto de detectores paralelos.

O segundo é o \textbf{drone como instrumento de verificação, e não apenas de
patrulha}. A inspeção aérea programada tem valor, mas o ganho maior está em
responder a um alerta específico levando um sensor móvel à coordenada em
segundos — o que transforma uma dúvida de vinte minutos em uma confirmação de
dois.

O terceiro é tratar a \textbf{degradação do próprio sistema como evento de
segurança}. Uma câmera cega não é um chamado de manutenção: é um ponto cego de
detecção, e deve chegar ao operador com a mesma clareza de um alerta de
intrusão. Essa inversão é incomum e barata de implementar.

O quarto é a \textbf{evidência como requisito de origem}. Por operar em
ambiente regulado, o sistema nasce com trilha imutável de ocorrências e ações,
em vez de ganhar um módulo de auditoria depois — decisão que a experiência do
grupo em artefatos anteriores mostrou ser muito mais cara quando postergada.

\subsection{Premissas, restrições e riscos}

A Tabela~\ref{tab:riscos} registra as condições sob as quais a proposta se
sustenta e os riscos que a equipe já identifica. Registrá-los na Sprint 1 é
parte do trabalho: proposta sem restrição declarada costuma ser proposta mal
compreendida.

\begin{table}[H]
\centering
\caption{Premissas, restrições e riscos identificados}
\label{tab:riscos}
{\small\singlespacing
\begin{tabularx}{\textwidth}{L{2.0cm} L{4.4cm} X}
\toprule
\cab{Tipo} & \cab{Descrição} & \cab{Tratamento previsto} \\
\midrule
Premissa & A instalação já dispõe de perímetro físico, controle de acesso e
CFTV & O ARGOS integra o que existe; não substitui barreira física \\
Premissa & Há equipe de resposta disponível em turno & Sem resposta, detecção
mais rápida não produz proteção \\
Restrição & Operação de drone sujeita à regulamentação aeronáutica: RBAC-E 100,
da ANAC, que substituiu o RBAC-E 94, e ICA 100-40, do DECEA & Rotas, altitudes
e autorizações tratadas como requisito, não como detalhe de operação; voos além
da linha de visada exigem requisitos mais rigorosos \\
Restrição & Tratamento de imagem de pessoas sujeito à LGPD — Lei n.º 13.709/2018
—, que classifica dado biométrico como dado pessoal sensível & Minimização,
retenção definida, controle de acesso por perfil e registro de quem acessou \\
Restrição & Ambiente de defesa restringe tráfego de dados para fora da
instalação & Processamento na borda e plataforma hospedada internamente \\
Restrição & Requisitos de proteção física definidos pela CNEN NN 2.01 & O
sistema apoia o Sistema de Proteção Física existente, sem substituí-lo \\
Risco & Falso positivo excessivo levar o operador a ignorar alertas & Correlação
entre fontes, política de severidade e realimentação pela classificação do
operador \\
Risco & Dependência de um único fornecedor de hardware & Interfaces abertas e
camada de adaptação por tipo de sensor \\
Risco & Expectativa de autonomia total da IA & O documento afirma, e o produto
reforça, que a decisão é humana e assistida \\
Risco & Indisponibilidade de dados reais para treinar e avaliar os modelos &
Uso inicial de conjuntos públicos e de dados sintéticos, com reavaliação quando
houver acesso \\
\bottomrule
\end{tabularx}}
\fonte{elaborado pelos autores (2026).}
\end{table}

\subsection{Registro de decisão arquitetural}
\label{sec:adr}

A Tabela~\ref{tab:adr} registra, no modelo de Nygard (2011), a primeira decisão
de consequência tomada pela equipe. Registrá-la agora evita que ela seja
revisitada por esquecimento nas Sprints seguintes.

\begin{table}[H]
\centering
\caption{ADR-001 — Processamento de visão computacional na borda}
\label{tab:adr}
{\small\singlespacing
\begin{tabularx}{\textwidth}{L{3.2cm} X}
\toprule
\cab{Identificador} & ADR-001 — Executar a detecção por visão computacional na
borda, junto à câmera \\
\midrule
Status & Aceita, na Sprint 1, como decisão preliminar de visão \\
Contexto & O sistema precisa detectar eventos em segundos; o ambiente é de
defesa, com restrição ao tráfego de imagens para fora da instalação; o enlace
entre pontos do perímetro pode ser instável; e o volume bruto de vídeo de
dezenas de câmeras é grande demais para ser transportado continuamente \\
Decisão & Executar a inferência de detecção em dispositivos de borda junto às
câmeras e aos drones, transmitindo à plataforma central apenas eventos
estruturados e os trechos de vídeo associados \\
Consequências positivas & Latência compatível com a meta de detecção abaixo de
60 s; menor consumo de enlace; menor exposição de imagem; detecção continua
operando se a plataforma central ficar indisponível \\
Consequências negativas & Maior custo de hardware por ponto; atualização de
modelo passa a ser um problema de distribuição em campo; capacidade de
processamento limitada restringe o tamanho dos modelos \\
Alternativas consideradas & Processamento centralizado na instalação,
descartado pela latência e pelo consumo de enlace; processamento em nuvem
pública, descartado pela restrição de tráfego de dados sensíveis em ambiente de
defesa \\
Gatilho de revisão & A decisão deve ser revisitada se a instalação passar a
dispor de rede interna de alta capacidade com baixa latência, ou se os modelos
exigidos ultrapassarem a capacidade dos dispositivos de borda \\
\bottomrule
\end{tabularx}}
\fonte{elaborado pelos autores (2026), conforme o modelo de Nygard (2011).}
\end{table}

% =========================================================== 6
\section{BENEFÍCIOS ESPERADOS}
\label{sec:beneficios}

Os benefícios foram organizados em três grupos — operacionais, de segurança e
de conformidade —, cada um associado a um indicador verificável. A
Tabela~\ref{tab:beneficios} evita a formulação genérica (``mais segurança'',
``mais eficiência'') em favor de afirmações que podem ser confirmadas ou
desmentidas por medição.

\begin{table}[H]
\centering
\caption{Benefícios esperados e como serão verificados}
\label{tab:beneficios}
{\small\singlespacing
\begin{tabularx}{\textwidth}{L{2.0cm} L{4.6cm} X}
\toprule
\cab{Grupo} & \cab{Benefício} & \cab{Como será verificado} \\
\midrule
Operacional & Atenção do operador dirigida ao que importa & Número de eventos
efetivamente julgados por turno e proporção de alertas descartados \\
Operacional & Verificação remota antes do deslocamento & Proporção de alertas
resolvidos sem envio de equipe e tempo médio até a confirmação visual \\
Operacional & Inspeção programada sem interromper a operação & Número de
inspeções executadas por período e horas de deslocamento evitadas \\
Segurança & Redução do tempo entre evento e compreensão & Tempo até a detecção,
comparado à linha de base após a implantação \\
Segurança & Cobertura contínua, inclusive onde não há câmera & Extensão do
perímetro com detecção automática ativa \\
Segurança & Ponto cego deixa de passar despercebido & Tempo médio entre a falha
de um sensor e o seu registro como evento \\
Conformidade & Evidência pronta para auditoria & Proporção de ocorrências com
trilha completa e tempo para produzir um relatório ao regulador \\
Conformidade & Tratamento de imagem com controle & Registro de acesso às
imagens e aderência à política de retenção \\
Institucional & Base replicável para outras instalações críticas & Esforço de
implantação em uma segunda instalação, comparado ao da primeira \\
\bottomrule
\end{tabularx}}
\fonte{elaborado pelos autores (2026).}
\end{table}

Um benefício merece registro em separado, por ser o menos visível. Ao obrigar o
operador a classificar cada alerta — confirmado ou falso positivo, e por qual
motivo —, o sistema produz, como subproduto da operação diária, um conjunto de
dados rotulados sobre o que de fato acontece naquela instalação. Esse acervo é
o que permite reduzir o falso positivo ao longo do tempo e, mais do que isso, é
conhecimento que hoje se perde na memória de quem estava no turno.

% =========================================================== 7
\section{ESCOPO}
\label{sec:escopo}

\subsection{Dentro do escopo}

A proposta da Sprint 1 compreende, como escopo do sistema:

\begin{itemize}[leftmargin=*,itemsep=2pt,topsep=4pt]
  \item monitorar o perímetro e as áreas restritas por meio de sensores,
        câmeras com visão computacional e drones;
  \item detectar intrusão, permanência indevida, objeto abandonado e violação
        de barreira;
  \item correlacionar fontes independentes e atribuir severidade ao evento;
  \item gerar alertas automáticos e ordená-los por risco na fila do operador;
  \item exibir a situação em painel operacional com mapa, vídeo e indicadores;
  \item despachar drone para verificação remota de um alerta e transmitir a
        imagem ao operador;
  \item executar rotas de inspeção programadas e registrar o resultado;
  \item registrar ocorrências, ações e evidências em trilha auditável;
  \item monitorar a saúde dos próprios sensores, câmeras e enlaces;
  \item administrar sensores, rotas, perfis de acesso e políticas de
        severidade.
\end{itemize}

\subsection{Fora do escopo}

Delimitar o que não será feito é parte do trabalho desta Sprint: proposta que
promete tudo não é ambiciosa, é imprecisa. A Tabela~\ref{tab:fora} registra as
exclusões e, principalmente, a razão de cada uma.

\begin{table}[H]
\centering
\caption{Itens deliberadamente fora do escopo}
\label{tab:fora}
{\small\singlespacing
\begin{tabularx}{\textwidth}{L{5.0cm} X}
\toprule
\cab{Fora do escopo} & \cab{Por quê} \\
\midrule
Monitoração de segurança nuclear do processo (\emph{safety}): reator,
refrigeração, barreiras de contenção & Domínio distinto, com sistemas
dedicados, redundantes e certificados; o ARGOS trata de proteção física
(\emph{security}) \\
Resposta automática sem intervenção humana — travar portas, acionar sirene,
despachar equipe sem confirmação & A decisão deve permanecer humana em ambiente
regulado; o sistema recomenda e registra, não decide \\
Identificação biométrica de pessoas e reconhecimento facial & Risco
desproporcional em relação ao benefício no recorte tratado, e tratamento de dado
pessoal sensível sob a LGPD; o sistema detecta presença humana, não identidade \\
Substituição do sistema de controle de acesso existente & O ARGOS consome a
informação de credenciais; não gerencia credenciamento \\
Contramedidas ativas contra drones intrusos (interferência, captura, abate) &
Atividade regulada e restrita a autoridade competente \\
Fabricação de hardware (drones, sensores, câmeras) & O projeto integra
equipamento de mercado por interfaces abertas \\
Operação em instalações não nucleares & Possível no futuro; manter o foco é o
que torna a Sprint executável \\
Aplicativo móvel completo para o público externo & Fora do recorte do centro de
operações nesta etapa \\
\bottomrule
\end{tabularx}}
\fonte{elaborado pelos autores (2026).}
\end{table}

\subsection{Requisitos não funcionais iniciais}

Os requisitos da Tabela~\ref{tab:rnf} serão detalhados nas próximas Sprints,
mas já condicionam decisões de arquitetura e por isso são declarados agora.

\begin{table}[H]
\centering
\caption{Requisitos não funcionais iniciais}
\label{tab:rnf}
{\small\singlespacing
\begin{tabularx}{\textwidth}{L{1.6cm} L{3.4cm} X}
\toprule
\cab{Código} & \cab{Atributo} & \cab{Requisito} \\
\midrule
RNF-01 & Desempenho & Do evento ao alerta no painel em menos de 60 s, no
percentil 95 \\
RNF-02 & Desempenho & Inferência na borda concluída em menos de 200 ms por
quadro analisado \\
RNF-03 & Disponibilidade & Disponibilidade mensal de 99,5 \% da função de
detecção \\
RNF-04 & Resiliência & A detecção na borda continua operando por no mínimo 4 h
sem conexão com a plataforma central \\
RNF-05 & Segurança & Autenticação forte, perfis de acesso e registro de toda
ação do operador \\
RNF-06 & Auditabilidade & Trilha de ocorrências imutável e exportável em
formato aberto \\
RNF-07 & Privacidade & Retenção definida por política, com minimização e
registro de acesso às imagens \\
RNF-08 & Interoperabilidade & Integração por interfaces abertas com CFTV,
controle de acesso e sensores de diferentes fabricantes \\
RNF-09 & Usabilidade & Alerta compreendido e julgado em até 30 s por operador
treinado, sem sair da tela de verificação \\
RNF-10 & Escalabilidade & Crescimento do parque de sensores sem reprojeto da
plataforma \\
RNF-11 & Conformidade & Aderência à CNEN NN 2.01, à LGPD e à regulamentação
aeronáutica aplicável à operação dos drones \\
\bottomrule
\end{tabularx}}
\fonte{elaborado pelos autores (2026).}
\end{table}

\subsection{Critérios de aceitação desta Sprint}

Em coerência com a prática de Definição de Pronto adotada pelo grupo na
disciplina, a Sprint 1 é considerada concluída quando: o problema estiver
descrito com as situações anômalas identificadas; os stakeholders estiverem
mapeados com papel e necessidade; a solução tiver nome, funcionalidades,
tecnologias e diferenciais declarados; o escopo estiver delimitado nos dois
sentidos; houver protótipo conceitual cobrindo tela principal, fluxo e
funcionalidades; e o documento estiver versionado no repositório do projeto.
Todos os itens foram cumpridos nesta entrega.

% =========================================================== 8
\section{PROTÓTIPO CONCEITUAL}
\label{sec:prototipo}

O protótipo é de baixa fidelidade e tem finalidade exclusivamente conceitual:
serve para verificar se o fluxo de trabalho proposto faz sentido para quem
operará o sistema, e não para definir acabamento visual, paleta ou tipografia.
As telas foram desenhadas a partir da pergunta que o operador faz em cada
momento, e não a partir das funcionalidades disponíveis.

\subsection{Fluxo geral de utilização}

A Figura~\ref{fig:fluxo} apresenta o fluxo completo, do evento captado até o
registro da ocorrência, incluindo o caminho do falso positivo — que não é um
desvio do processo, e sim parte dele: é o que realimenta o modelo.

\begin{figure}[H]
\centering
\caption{Fluxo geral de utilização do ARGOS}
\label{fig:fluxo}
\includegraphics[width=\textwidth]{fig_fluxo.png}
\fonte{elaborado pelos autores (2026).}
\end{figure}

\subsection{Tela principal — painel operacional}

A tela da Figura~\ref{fig:painel} responde à pergunta ``o que está acontecendo
agora?''. Ela reúne o mapa da instalação com a posição dos sensores, das
câmeras e dos drones, a fila de alertas ordenada por severidade, os indicadores
do turno e o histórico recente. Cada alerta traz, desde a fila, as duas ações
mais prováveis — verificar e despachar —, para que o operador não precise
navegar antes de agir.

\begin{figure}[H]
\centering
\caption{Protótipo conceitual — painel operacional}
\label{fig:painel}
\includegraphics[width=\textwidth]{fig_tela_painel.png}
\fonte{elaborado pelos autores (2026).}
\end{figure}

\subsection{Tela de verificação da ocorrência}

A tela da Figura~\ref{fig:ocorrencia} responde a ``o que é isto, exatamente?''.
Reúne em uma única visão o quadro de vídeo com as detecções do modelo e seus
graus de confiança, a linha do tempo do evento, o contexto que sustenta a
classificação de risco — inclusive a ausência de credencial registrada no setor
— e as quatro ações possíveis. A exibição do grau de confiança é deliberada: o
operador precisa saber o quanto o modelo está certo para decidir quanto peso
dar a ele.

\begin{figure}[H]
\centering
\caption{Protótipo conceitual — verificação da ocorrência}
\label{fig:ocorrencia}
\includegraphics[width=\textwidth]{fig_tela_ocorrencia.png}
\fonte{elaborado pelos autores (2026).}
\end{figure}

\subsection{Tela de inspeção aérea}

A tela da Figura~\ref{fig:inspecao} responde a ``o que o drone está vendo?''.
Apresenta a rota com os pontos de passagem, a posição atual da aeronave, a
telemetria necessária para decidir se a missão pode continuar e a transmissão
ao vivo. O \emph{checklist} da missão aparece na própria tela porque a
liberação do espaço aéreo é condição de partida, e não formalidade: a
operação está sujeita à regulamentação da ANAC e do DECEA, conforme registrado
na Tabela~\ref{tab:riscos}.

\begin{figure}[H]
\centering
\caption{Protótipo conceitual — inspeção aérea por drone}
\label{fig:inspecao}
\includegraphics[width=\textwidth]{fig_tela_inspecao.png}
\fonte{elaborado pelos autores (2026).}
\end{figure}

\subsection{Observações sobre o protótipo}

Três escolhas do protótipo expressam decisões de projeto e não preferências
estéticas. A primeira: a fila de alertas é curta e ordenada por severidade,
porque uma lista longa reproduz, na tela, o mesmo problema de atenção que o
sistema pretende resolver. A segunda: o botão de classificar como falso
positivo tem a mesma proeminência dos demais, porque o registro do erro é tão
valioso quanto o do acerto. A terceira: nenhuma tela executa ação irreversível
sem confirmação explícita do operador, coerentemente com a decisão de manter a
autonomia do sistema limitada à recomendação.

% =========================================================== 9
\section{CONCLUSÃO}
\label{sec:conclusao}

A Sprint 1 cumpriu o que se propôs: compreender o problema antes de desenhar a
solução. O exercício mostrou que o desafio proposto pela AMAZUL, apesar de
enunciado em termos de tecnologias — drones, sensores, visão computacional,
Inteligência Artificial —, é, no fundo, um problema de tempo e de atenção. A
proteção física só é efetiva se a detecção somada à resposta couber dentro do
tempo que as barreiras físicas concedem, e o elo frágil dessa cadeia, hoje, é a
detecção dependente de atenção humana sustentada.

Daí decorre a proposta. O ARGOS não é um conjunto de câmeras inteligentes: é
uma camada de correlação que transforma detecções isoladas em eventos
classificados por risco, acrescenta ao operador a capacidade de verificar
remotamente antes de deslocar pessoas e preserva de cada decisão uma trilha
auditável. As três decisões mais importantes da Sprint foram, por isso,
decisões de limite: manter a decisão final sob responsabilidade humana, excluir
a identificação biométrica do escopo e tratar a degradação do próprio sistema
como evento de segurança.

Registre-se também o que esta Sprint \emph{não} produziu. Não há arquitetura
decidida, não há modelo de dados, não há avaliação de viabilidade de hardware e
não há validação com usuários reais — e os números de linha de base são
estimativas da equipe, não medições. Apresentar qualquer uma dessas lacunas
como resolvida comprometeria o valor do documento como base para as etapas
seguintes.

Os próximos passos previstos são: detalhar os casos de uso com fluxos principal
e alternativos; produzir o modelo de domínio e os diagramas de contêineres e de
componentes; avaliar a viabilidade técnica dos modelos de visão computacional
na borda; especificar as interfaces de integração com CFTV, controle de acesso
e sensores; e validar o protótipo com usuários que exerçam, na prática, o papel
de operador de monitoramento. A arquitetura, a modelagem e o plano de testes
decorrentes serão objeto das Sprints seguintes.

% =========================================================== REFERÊNCIAS
\newpage
\phantomsection
\addcontentsline{toc}{section}{REFERÊNCIAS}
\section*{REFERÊNCIAS}
\vspace{0.6\baselineskip}

{\singlespacing\raggedright
\setlength{\parindent}{0pt}
\setlength{\parskip}{0.7\baselineskip}

AGÊNCIA NACIONAL DE AVIAÇÃO CIVIL. \textbf{RBAC-E n.º 100}: regulamento
brasileiro da aviação civil especial — aeronaves não tripuladas. Brasília:
ANAC, 2026. Regulamento que substitui o RBAC-E n.º 94 e adota classificação por
risco operacional.

ASSOCIAÇÃO BRASILEIRA DE NORMAS TÉCNICAS. \textbf{NBR 14724}: informação e
documentação: trabalhos acadêmicos: apresentação. Rio de Janeiro, 2011.

ASSOCIAÇÃO BRASILEIRA DE NORMAS TÉCNICAS. \textbf{NBR 6023}: informação e
documentação: referências: elaboração. Rio de Janeiro, 2018.

BRASIL. \textbf{Lei n.º 12.706, de 8 de agosto de 2012}. Autoriza o Poder
Executivo a criar a empresa pública Amazônia Azul Tecnologias de Defesa S.A. —
AMAZUL. Brasília: Presidência da República, 2012. Disponível em:
https://www.planalto.gov.br/ccivil\_03/\_ato2011-2014/2012/lei/L12706.htm.
Acesso em: 30 out.\ 2026.

BRASIL. \textbf{Lei n.º 13.709, de 14 de agosto de 2018}. Lei Geral de Proteção
de Dados Pessoais (LGPD). Brasília: Presidência da República, 2018.

COMANDO DA AERONÁUTICA. \textbf{ICA 100-40}: aeronaves não tripuladas e o
acesso ao espaço aéreo brasileiro. Rio de Janeiro: DECEA, 2026.

COMISSÃO NACIONAL DE ENERGIA NUCLEAR. \textbf{Norma CNEN NN 2.01}: proteção
física de materiais e instalações nucleares. Rio de Janeiro: CNEN, 2019.
Aprovada pela Resolução CNEN n.º 253, de 11 de novembro de 2019.

INTERNATIONAL ATOMIC ENERGY AGENCY. \textbf{Nuclear Security Recommendations on
Physical Protection of Nuclear Material and Nuclear Facilities}
(INFCIRC/225/Revision 5). Vienna: IAEA, 2011. (IAEA Nuclear Security Series
n.º 13).

NYGARD, Michael. \textbf{Documenting architecture decisions}. 2011. Cognitect
Blog. Disponível em:
https://www.cognitect.com/blog/2011/11/15/documenting-architecture-decisions.
Acesso em: 30 out.\ 2026.

OBJECT MANAGEMENT GROUP. \textbf{OMG Unified Modeling Language (OMG UML)},
version 2.5.1. 2017. Especificação técnica formal/2017-12-05.

SILVA, Leandro Ramos da. \textbf{Software Engineering}: material de aula. São
Paulo: FIAP, 2026.

}

% =========================================================== APÊNDICE
\newpage
\phantomsection
\addcontentsline{toc}{section}{APÊNDICE A — MATRIZ DE RASTREABILIDADE}
\section*{APÊNDICE A — MATRIZ DE RASTREABILIDADE}
\vspace{0.6\baselineskip}

A matriz da Tabela~\ref{tab:rastreabilidade} liga cada necessidade identificada
ao stakeholder que a expressa, às funcionalidades que a atendem e ao benefício
esperado. Ela serve de verificação interna: necessidade sem requisito indica
lacuna; requisito sem necessidade indica funcionalidade sem dono.

\begin{table}[H]
\centering
\caption{Rastreabilidade entre necessidades, requisitos e benefícios}
\label{tab:rastreabilidade}
{\small\singlespacing
\begin{tabularx}{\textwidth}{X L{3.0cm} L{2.6cm} L{3.4cm}}
\toprule
\cab{Necessidade} & \cab{Stakeholder} & \cab{Requisitos} & \cab{Benefício
esperado} \\
\midrule
Perceber a intrusão enquanto ela acontece & Operador; equipe de segurança &
RF-01, RF-02, RF-03 & Redução do tempo entre evento e compreensão \\
Não ser inundado por alarme falso & Operador & RF-03, RF-04, RF-13 & Atenção
dirigida ao que importa \\
Saber o que esperar antes de chegar ao local & Equipe de segurança física &
RF-06, RF-08 & Verificação remota antes do deslocamento \\
Inspecionar sem interromper a operação & Equipe de inspeção & RF-09 & Inspeção
programada com histórico comparável \\
Acompanhar o risco de forma consolidada & Gestor da instalação & RF-05, RF-11 &
Indicadores confiáveis para a direção \\
Dimensionar turnos pela carga real & Coordenador do centro de operações &
RF-11 & Distribuição de efetivo baseada em dado \\
Configurar sem depender do fornecedor & Administrador do sistema & RF-12 &
Autonomia operacional da equipe \\
Comprovar conformidade com evidência íntegra & Órgãos reguladores; gestor &
RF-07, RF-11 & Evidência pronta para auditoria \\
Não expor imagem nem dado sensível & TI e segurança da informação & RF-07,
RNF-05, RNF-07 & Tratamento de imagem com controle \\
Descobrir rapidamente o ponto cego & Operador; administrador & RF-10 & Falha de
sensor deixa de passar despercebida \\
Não depender de um único fabricante & Fornecedores; TI & RNF-08 & Base
replicável para outras instalações \\
\bottomrule
\end{tabularx}}
\fonte{elaborado pelos autores (2026).}
\end{table}

\end{document}
